%!TEX root = ../main/main.tex

To model extraction under errors, we consider documents that may differ from an ideal match by a bounded number of edit operations. In this section we define the edit model used throughout the paper. We first define edit operations at the level of documents, and then define their canonical action on spans and mappings.

\subsection{Edit operations}

Let $d=a_1\cdots a_n\in\Sigma^*$ be a document. For $p\in\mathbb{N}$ and $c\in\Sigma$, we consider the following atomic edit operations.

An \emph{insertion} $\ins_{p,c}$ is valid on $d$ if $1\leq p\leq n+1$. Its effect is
\[
\ins_{p,c}(d)
=
a_1\cdots a_{p-1}\,c\,a_p\cdots a_n.
\]
In particular, the case $p=n+1$ appends $c$ at the end of the document.

A \emph{deletion} $\del_p$ is valid on $d$ if $1\leq p\leq n$. Its effect is
\[
\del_p(d)
=
a_1\cdots a_{p-1}a_{p+1}\cdots a_n.
\]

A \emph{substitution} $\sub_{p,c}$ is valid on $d$ if $1\leq p\leq n$. Its effect is
\[
\sub_{p,c}(d)
=
a_1\cdots a_{p-1}\,c\,a_{p+1}\cdots a_n.
\]


Let
\[
\operatorname{Edits}
=
\{\ins_{p,c}\mid p\in\mathbb{N},\,c\in\Sigma\}
\cup
\{\del_p\mid p\in\mathbb{N}\}
\cup
\{\sub_{p,c}\mid p\in\mathbb{N},\,c\in\Sigma\}
\]
be the set of all atomic edit operations over $\Sigma$.

\subsection{Edit scripts}

An \emph{edit script} is a finite sequence
\[
E=e_1\cdots e_m\in\operatorname{Edits}^*.
\]
The empty edit script is denoted by $\epsilon$. A script $E=e_1\cdots e_m$ is \emph{valid} on a document $d$ if $e_1$ is valid on $d$, $e_2$ is valid on $e_1(d)$, and so on.

The action of a valid edit script on $d$ is defined by
\[
\epsilon(d)=d
\]
and
\[
E(d)=e_m(\cdots e_2(e_1(d))\cdots).
\]
Equivalently, the operations of $E$ are applied from left to right. We write $|E|=m$ for the length of the script.

The \emph{edit distance} between two documents $d,d'\in\Sigma^*$ is denoted by
\[
\operatorname{edist}(d,d')
\]
and is defined as
\[
\operatorname{edist}(d,d')
=
\min
\{
|E|
\mid
E\in\operatorname{Edits}^*,
\ E \text{ is valid on } d,
\ E(d)=d'
\}.
\]

\begin{example}
	Let $d=\mathtt{a}$ and $d'=\mathtt{abc}$. Then
	\[
	E=\ins_{2,\mathtt{b}}\,\ins_{3,\mathtt{c}}
	\]
	is a valid edit script on $d$ such that $E(d)=d'$. Indeed,
	\[
	\mathtt{a}
	\xrightarrow{\ins_{2,\mathtt{b}}}
	\mathtt{ab}
	\xrightarrow{\ins_{3,\mathtt{c}}}
	\mathtt{abc}.
	\]
	At least two insertions are necessary therefore
	\[
	\operatorname{edist}(\mathtt{a},\mathtt{abc})=2.
	\]
\end{example}


\santiago{Lo anterior deberia estar bien. Lo que viene a continuacion es todo lo nuevo a revisar}

\subsection{Edit distance in information extraction}

Given a document \(d\) and a query
\(S\), an approximate extraction semantics should explain when a mapping
\(\nu\) over \(d\) is justified by the existence of an edited document \(d'\)
such that \(S\) matches \(d'\) exactly.

More concretely, suppose that an edit script \(E\) transforms \(d\) into \(d'\),
and that
\[
\mu\in\sem{S}{d'}.
\]
Then the central question is how the mapping \(\mu\), whose spans refer to
positions in \(d'\), induces an output mapping \(\nu\), whose spans must refer to
positions in the original document \(d\).

A first, natural idea is the following. Let \(E\) be an edit script such that
\(E(d)=d'\), and let
\[
\mu\in\sem{S}{d'}.
\]
Then one would like to transport \(\mu\) back through the edit script \(E\), in
order to obtain a mapping over \(d\). This suggests a semantics of the form
\[
\nu = E^{-1}(\mu).
\]
Under this view, approximate extraction would return all mappings \(\nu\) for
which there exists an edited document \(d'\), an edit script \(E\) from \(d\) to
\(d'\), and an exact match \(\mu\in\sem{S}{d'}\), such that \(\nu\) is obtained
by transporting \(\mu\) back along \(E\).

%
%\begin{example}	
%	Consider the document $d=\mathtt{ac}.$	
%	First, consider the query
%	\[
%	S_1=\mathtt a\cdot \mathtt b\cdot x\{\mathtt c\}.
%	\]
%	By inserting \(\mathtt b\) at position \(2\), we obtain $d'=\mathtt{abc}$. On \(d'\), the exact match assigns
%	\[
%	\mu_1(x)=\Span{3,4}.
%	\]
%	Transporting this match back to the original document \(d\), the inserted
%	symbol \(\mathtt b\) disappears. Since this symbol occurs before the opening
%	of \(x\), the resulting output should be
%	\[
%	\nu_1(x)=\Span{2,3}.
%	\]
%	
%This result means that we capture de letter $c$ in the %original document $d$ but we have to insert a letter $b$ to match the query $S_1$.

%Now consider instead the query
%\[
%S_2=\mathtt a\cdot x\{\mathtt b\mathtt c\}.
%\]
%Using the same document edit as before, we obtain 
%$ d=\mathtt{ac} \xrightarrow{\ins_{2,\mathtt b}} d'=\mathtt{abc},$ the exact match on \(d'\) assigns
%\[
%\mu_2(x)=\Span{2,4}.
%\]
%In this case, the inserted symbol \(\mathtt b\) belongs to the capture of
%\(x\). After transporting the match back to \(d\), we again should obtain
%\[
%\nu_2(x)=\Span{2,3}.
%\]

%Thus, both cases use the same document edit and produce the same output
%mapping over the original document:
%\[
%\nu_1(x)=\nu_2(x)=\Span{2,3}.
%\]
%However, the inserted symbol is outside the capture in the first case and
%inside the capture in the second one. Therefore, the plain edit operation
%\(\ins_{2,\mathtt b}\) does not determine its effect on captured spans.
%\end{example}


\subsection{Plain edit scripts are not enough}

Let see an example to introduce the problem:

\begin{example}
	Consider the query $S=\mathtt a\cdot x\{\mathtt c\}$
	over the document $d=\mathtt{abc}.$
	By deleting the symbol \(\mathtt b\) at position \(2\), we obtain
	\[
	d'=\mathtt{ac}.
	\]
	On \(d'\), the exact match assigns
	\[
	\mu(x)=\Span{2,3}.
	\]
	
	Now suppose that we want to transport this exact match back to the original
	document \(d\). Since \(d'\) was obtained from \(d\) by deleting
	\(\mathtt b\), transporting the mapping back amounts to inserting
	\(\mathtt b\) again at position \(2\).
	
	However, this insertion can be interpreted in two different ways with
	respect to the capture of \(x\).
	
	First, the deleted symbol may have occurred before opening \(x\). In that
	case, the run is intuitively
	\[
	\mathtt a\;\; \del(\mathtt b)\;\; x\{\mathtt c\},
	\]
	and after transporting the match back to \(d\), the output should be
	\[
	\nu_1(x)=\Span{3,4}.
	\]
	This means that \(x\) captures only the symbol \(\mathtt c\) in the original
	document.
	
	Second, the deleted symbol may have occurred after opening \(x\). In that
	case, the run is intuitively
	\[
	\mathtt a\;\; x\{\del(\mathtt b)\;\mathtt c\},
	\]
	and after transporting the match back to \(d\), the output should be
	\[
	\nu_2(x)=\Span{2,4}.
	\]
	This means that \(x\) captures the substring \(\mathtt{bc}\) in the original
	document.
	
\end{example}

The problem is that an inverse delete operation does not determine how the insertion should interact with the variable markers. In the last example, the symbol \(\mathtt b\) may be reinserted
before the capture of \(x\), or inside the capture of \(x\), leading to
different output mappings and both ways are compatible with the mapping over de edited document.


For this reason, in our definition of a approximate extraction, edit operations are allowed to occur at different points of an extraction run and their effect on the output mapping is determined by their position relative to the opening and closing markers of variables.


\subsection{The solution}

\subsection{Approximate variable-set automata}



Let $\mathcal{A}=(Q,\Sigma,V,q_0,F,\delta)$
be a variable-set automaton, and let \(k\in\mathbb{N}\). We define the
\(k\)-approximate automaton of \(\mathcal{A}\) as
\[
\mathcal{A}_{\leq k}=(Q',\Sigma,V,q'_0,F',\delta'),
\]
where
\begin{itemize}
	\item $Q'=Q\times\{0,\ldots,k\}$
	\item $q'_0=(q_0,0)$
	\item $F'=F\times\{0,\ldots,k\}$
\end{itemize}


The second component of a state counts the number of edit operations used so
far.

The transition relation \(\delta'\) is defined as follows.

\paragraph{Variable transitions.}
For every variable marker transition
\[
(q,\gamma,q')\in\delta,
\]
where \(\gamma\in\{\ox,\cx\mid x\in V\}\), and every \(r\leq k\), we add
\[
((q,r),\gamma,(q',r))\in\delta'.
\]
Variable transitions do not consume input and do not increase the edit cost.


\paragraph{Exact letter transitions.}
For every letter transition
\[
(q,c,q')\in\delta
\]
and every \(r\leq k\), we add
\[
((q,r),c,(q',r))\in\delta'.
\]
This corresponds to reading the expected symbol without using an edit operation.

\paragraph{Substitution transitions.}
For every letter transition
\[
(q,c,q')\in\delta,
\]
every \(a\in\Sigma\) with \(a\neq c\), and every \(r<k\), we add
\[
((q,r),a,(q',r+1))\in\delta'.
\]
This consumes the symbol \(a\) from the original document while simulating the
symbol \(c\) expected by \(\mathcal{A}\).

\paragraph{Insertion transitions.}
For every letter transition
\[
(q,c,q')\in\delta
\]
and every \(r<k\), we add
\[
((q,r),\varepsilon,(q',r+1))\in\delta'.
\]
This simulates that the symbol \(c\) is inserted in the edited document. The
transition advances in the original automaton but consumes no symbol from the
original document.

\paragraph{Deletion transitions.}
For every state \(q\in Q\), every \(a\in\Sigma\), and every \(r<k\), we add
\[
((q,r),a,(q,r+1))\in\delta'.
\]
This consumes the symbol \(a\) from the original document without advancing in
the original automaton. Intuitively, the symbol is deleted before the exact
match is performed.

The semantics of \(\mathcal{A}_{\leq k}\) is the ordinary semantics of a
variable-set automaton over the original input document. That is, for a document
\(d\), every accepting run of \(\mathcal{A}_{\leq k}\) over \(d\) produces a
mapping whose spans are boundaries of \(d\).

The important point is that different placements of the same edit operation are
represented by different runs. For instance, a deletion transition taken before
the opening marker \(x\vdash\) deletes a symbol outside the capture of \(x\),
whereas the same deletion transition taken after \(x\vdash\) and before
\(\dashv x\) deletes a symbol inside the capture of \(x\). Thus, the position of
an edit relative to the variable markers is determined by the run itself.